\section{Machine-learning strategy for autonomous damage detection}\label{sec:ml}

The machine-learning (ML) layer is the component that raises the proposed system from conventional Structural Health Monitoring (SHM), which typically terminates at data acquisition and storage, to genuine Structural Cognition, in which the structure interprets its own responses. This layer is organized as a statistical pattern-recognition pipeline whose four canonical procedures are operational evaluation, data acquisition, feature extraction, and statistical modelling for feature discrimination \citep{farrar2013shm,farrar2007introduction}. Fig.~\ref{fig:mlpipeline} shows how this pipeline maps onto the Rytter damage stages, from existence through prognosis. Within that pipeline, the central methodological commitment of the present framework is that ML does not operate on raw signals alone but on damage-sensitive features and on residuals computed against the calibrated digital twin, so that physics and data are fused rather than opposed \citep{brunton2019data}. The strategy is deliberately tiered, because the single most important constraint in any field SHM deployment is data scarcity: damaged-state data from a high-value structure are rarely available, so the architecture must function in a predominantly data-poor regime while remaining able to exploit labelled scenarios when they exist \citep{farrar2013shm,naser2023ml,avci2021review}.

\subsection{Damage-sensitive feature extraction}\label{subsec:features}

Feature extraction converts sensor data into low-dimensional quantities that are highly sensitive to structural condition and, ideally, insensitive to benign variability. Two complementary feature families are extracted. The first is a set of physics-based modal features: natural frequencies, mode shapes, and modal damping ratios, identified from the triaxial-accelerometer network under the rotating-unbalance excitation. These are physically interpretable because, for the idealized braced frame, the natural frequencies and modes depend only on the mass and stiffness distribution, so any stiffness loss, bracing degradation, or bolted-connection loosening manifests as a measurable downward shift in frequency, a distortion of the mode shape, and, where energy dissipation increases, a rise in measured damping \citep{chopra2017dynamics,farrar2013shm}. The second family is data-driven features obtained by compressing the distributed strain and acceleration streams: the singular value decomposition / principal component analysis yields the dominant spatial modes of healthy operation, and dynamic mode decomposition extracts spatio-temporal coherent structures, each with an associated oscillation frequency and growth/decay rate, directly from experimental snapshots without requiring the governing equations \citep{brunton2019data}. A third, and arguably the most discriminative, feature is the residual between the measured response and the digital-twin prediction. Because the finite-element-based twin produces baseline modal parameters and time histories, prediction errors and statistical indicators on these residuals act as features whose growth flags a divergence between the physical structure and its validated model \citep{brunton2019data,farrar2013shm}.

\subsection{Environmental and operational normalization}\label{subsec:normalization}

A feature is useful for damage detection only after benign operational and environmental variability has been separated from genuine condition change; otherwise temperature-driven stiffness variation or changes in excitation amplitude are mistaken for damage, producing false-positive alarms \citep{farrar2013shm,figueiredo2010machine}. The planned temperature sensors provide the auxiliary measurements needed for this normalization, supplying thermal compensation of the strain channels and allowing modal features to be conditioned on temperature. Temperature is the confounder most directly measured and compensated, but it is not the only one. Humidity, wind buffeting and randomly varying operational loads also shift natural frequencies and damping; because the testbed instruments temperature rather than these quantities, their effect is handled statistically, by comparing features at similar points of the operational cycle and conditioning the healthy model on the inferred operational and environmental state \citep{figueiredo2010machine}. Normalization is treated as both a physical correction (thermal compensation of raw signals) and a statistical one (comparing features measured at similar points of the operational cycle, and conditioning the healthy model on the measured environmental state). Channel-specific care is required during fusion: tree ensembles are largely insensitive to scaling, whereas distance- and gradient-based learners (SVM, neural networks, PCA-derived features) are sensitive and require consistent normalization across the heterogeneous strain, acceleration, temperature, and displacement channels \citep{naser2023ml}.

\subsection{Unsupervised novelty detection (the data-poor default)}\label{subsec:unsupervised}

Because only healthy-state data are generally available, the default operating mode is unsupervised novelty/outlier detection: a model of the single healthy class is constructed, and incoming data are tested for consistency with that class \citep{farrar2013shm,wangcha2022unsupervised}. The primary tool is an autoencoder trained exclusively on healthy multi-sensor data; it learns a compact latent representation of normal behaviour, and its reconstruction error provides a natural, continuously evaluated anomaly indicator that rises when the structure departs from the learned baseline \citep{brunton2019data}. A linear autoencoder reduces to the SVD/PCA subspace, which makes the method interpretable and links the deep and classical formulations \citep{brunton2019data}. Complementary classical detectors provide low-parameter alternatives suited to streaming deployment: sequential probability ratio tests accumulate evidence across successive samples before declaring a change \citep{sohn2003statistical}, while outlier analysis based on Mahalanobis distance and density- or isolation-based methods offers distribution- and density-driven alternatives \citep{naser2023ml}. This route answers Rytter stage~1 (existence) and, when the spatially distributed sensor array localizes the anomalous response, contributes to stage~2 (location); it is subject to the axiomatic trade-off between damage sensitivity and noise-rejection capability, which governs threshold selection \citep{farrar2013shm}.

\subsection{Supervised classification (when labelled or simulated-damage data exist)}\label{subsec:supervised}

When damage classes can be labelled, most realistically from simulated-damage scenarios rather than from physical destruction of the testbed, supervised classifiers add localization and severity estimation. Random Forest, an ensemble of decision trees aggregated by majority voting, is non-parametric, insensitive to feature scaling, and pairs naturally with feature-importance analysis; Support Vector Machines separate damage classes by margin-maximizing hyperplanes with linear or radial-basis kernels; and Artificial Neural Networks, trained by back-propagation, capture more complex nonlinear mappings between features and damage state \citep{naser2023ml}. These supervised methods address Rytter stages~3 and~4 (type and severity), classifying conditions along an ordinal scale such as healthy / minor / major damage, consistent with the SHM axiom that existence and location are reachable unsupervised but type and severity generally require supervised learning \citep{farrar2013shm,naser2023ml,parisi2022steel}. Deployed studies confirm the maturity of this route for steel and infrastructure assets: supervised deep networks have classified the condition of bridge steel bearings \citep{wangsu2023steelbearing} and have detected and categorised damage across large, heterogeneous field datasets \citep{arya2021roaddamage}, illustrating the scale at which the supervised branch can operate once labelled data are available.

\subsection{Temporal modelling for progressive damage and bolt-loosening}\label{subsec:temporal}

Bolted-connection loosening and progressive stiffness loss are inherently time-evolving, and they are the damage mechanisms most directly relevant to the bolted SHS braced frame; loosening alters connectivity and adds energy dissipation, observable as drift in frequency and damping over time \citep{farrar2013shm}. Long Short-Term Memory (LSTM) recurrent networks are therefore used to model the temporal structure of the continuously acquired strain and vibration series, learning the healthy temporal evolution and forecasting expected response so that slow deviations, rather than single-sample outliers, are detected and trended \citep{brunton2019data,naser2023ml}. This supports incipient-fault detection and feeds the prognosis-oriented Rytter stage~5.

\subsection{Training, validation, evaluation, and explainability}\label{subsec:training}

Data scarcity is addressed by exploiting the validated finite-element model as a controlled data generator: parametric simulation of damage scenarios (localized stiffness reductions, bracing degradation, connection-flexibility changes) produces labelled cases that are compiled into a hybrid real-plus-simulated dataset for supervised training, with the explicit caveat that experimental data remain preferred and synthetic data require validation \citep{naser2023ml,cook2002fea}. Accordingly, any preliminary results obtained on this finite-element-generated, self-labelled set are reported as a specification of the validation protocol and its target metrics rather than as evidence of field performance, which only physically induced damage can establish. To avoid leakage between temporally adjacent windows of the same test, validation uses grouped (blocked) partitioning, in which all windows from a given test, damage scenario and operating or environmental condition are kept together in a single fold; generalisation is then assessed on an external test set of unseen scenarios and conditions rather than on randomly shuffled samples, and overfitting or underfitting is monitored across these grouped folds \citep{naser2023ml}. Performance is reported with several independent indicators rather than a single score, including the confusion matrix, ROC, precision, recall, F1, and balanced accuracy for classification, together with error and residual metrics for regression, while statistical models are explicitly tuned to balance false-positive and false-negative diagnoses according to their differing safety and economic costs \citep{farrar2013shm,naser2023ml}. To make autonomous integrity assessment trustworthy, decisions are accompanied by explainability, using permutation feature importance and SHAP, so that the features driving a damage indication can be checked for physical plausibility against the modal and residual quantities defined above \citep{naser2023ml}. The complete technique-to-stage mapping, summarized in Table~\ref{tab:mlmethods}, runs as follows: SVD/PCA, DMD, and twin-residual features supply the cognitive substrate; autoencoder and outlier analysis realize unsupervised perception (existence and location); Random Forest, SVM, and ANN realize supervised interpretation (type and severity); and LSTM realizes temporal reasoning toward prognosis. Together these methods advance the structure through successive Structural-Cognition levels rather than leaving it at passive data acquisition \citep{farrar2013shm,brunton2019data,naser2023ml}.
